\documentclass[11pt, letterpaper]{article}
\usepackage[utf8]{inputenc}
\usepackage[spanish, es-tabla]{babel}
\usepackage{graphicx} 
\usepackage{geometry}
\usepackage{amsmath, amssymb}
\geometry{left=2.5cm, right=2.5cm, top=2cm, bottom=2.5cm}
\usepackage{caption}
\usepackage{float}
\usepackage{url}
\usepackage{mathptmx}
\usepackage{titlesec}
\usepackage{microtype}
\usepackage{placeins}
\usepackage[T1]{fontenc}
\usepackage{array}
\usepackage[hidelinks]{hyperref}
\usepackage{tabularx}
\setlength{\parskip}{3pt}
\emergencystretch=2em

\titleformat{\section}{\normalfont\fontsize{16}{19}\bfseries}{\thesection}{1em}{}
\titleformat{\subsection}{\normalfont\fontsize{14}{17}\bfseries}{\thesubsection}{1em}{}
\titleformat{\subsubsection}{\normalfont\fontsize{12}{14.5}\bfseries}{\thesubsubsection}{1em}{}

\begin{document}

\begin{titlepage}
    \centering
    
    \begin{minipage}{0.15\textwidth}
         \includegraphics[width=\linewidth]{img/logo_ipn.png}
    \end{minipage}
    \hfill
    \begin{minipage}{0.65\textwidth}
        \centering
        \textbf{\Large INSTITUTO POLITÉCNICO NACIONAL}\\[0.2cm]
        \textbf{\large ESCUELA SUPERIOR DE CÓMPUTO}\\[0.2cm]
        \large DEPARTAMENTO DE INGENIERÍA EN SISTEMAS COMPUTACIONALES
    \end{minipage}
    \hfill
    \begin{minipage}{0.15\textwidth}
         \includegraphics[width=\linewidth]{img/logo_escom.png} 
    \end{minipage}

    \vspace{2cm}

    {\Large \textbf{Unidad de Aprendizaje:}} \\[0.5cm]
    {\Huge \textbf{Sistemas Operativos}}
    
    \vspace{1.5cm}

    {\Large \textbf{Práctica}} \\[0.5cm]
    {\Huge \textbf{Inicio de sesión e interfaces basadas en texto}}

    \vspace{1cm}

    \begin{flushright}
        \large
        \vspace{1.5cm}
        
        \textbf{Equipo:} \\
        Alvarez Garcia Martin \\
        Hernandez Castillo Emir Adrian \\
        Joaquín Rojas González \\
        Nava Alvarado Donovan Emilio \\[1cm]
        
        \textbf{Grupo:} \\
        4CV4 \\[1cm]
        
        \textbf{Profesor:} \\
        José Alfredo Jiménez Benítez
    \end{flushright}

    \vfill

    {\large Ciudad de México, 11 de octubre de 2026}

\end{titlepage}

\renewcommand{\contentsname}{Contenido}

% RESUMEN
\section*{Resumen}
\phantomsection
\addcontentsline{toc}{section}{Resumen}
El presente proyecto consiste en el desarrollo de un sistema de inicio de sesión e interfaces basadas en texto, implementado en lenguaje C mediante la biblioteca \texttt{ncurses}.
El sistema, denominado JDM E-TICKETS, permite realizar el registro de usuarios y la autenticación mediante una interfaz interactiva en terminal.
Su diseño integra mecanismos de validación de datos, procesamiento criptográfico de contraseñas y almacenamiento local de información, con el propósito de establecer un flujo de acceso que contemple tanto la interacción con el usuario como el manejo de sus credenciales.

La interfaz permite desplazarse entre las opciones disponibles mediante el teclado, utilizando las teclas de dirección, \texttt{ENTER} y \texttt{ESC}.
A través de los formularios se capturan los datos necesarios para crear una cuenta o iniciar sesión.
Durante el registro se verifican condiciones como la longitud mínima del nombre de usuario, la estructura de la dirección de correo electrónico, la disponibilidad de ambos identificadores, la complejidad de la contraseña y la coincidencia entre la contraseña y su confirmación.
Estas validaciones permiten controlar la información antes de completar el registro y proporcionan una respuesta cuando los datos ingresados no satisfacen los requisitos establecidos.

En materia de seguridad, el sistema utiliza la biblioteca \texttt{libsodium} para generar y verificar hashes de contraseñas mediante el algoritmo Argon2id.
De esta manera, las contraseñas originales no necesitan almacenarse directamente en el archivo de usuarios.
La información persistente se conserva localmente en \texttt{usuarios.txt}, donde cada registro contiene el nombre de usuario, el correo electrónico y el hash correspondiente.
Este mecanismo permite mantener las cuentas disponibles entre distintas ejecuciones de la aplicación, sin depender de un sistema gestor de bases de datos o de un servicio remoto.

La organización del código fuente separa el flujo principal, las funciones de autenticación y los componentes de la interfaz, facilitando la identificación de las responsabilidades de cada módulo.
Asimismo, la compilación se administra mediante un archivo \texttt{Makefile}, que define las opciones y las bibliotecas necesarias para generar el ejecutable.

El proyecto integra conceptos de programación en C, manejo de terminal, validación de entradas, persistencia en archivos y protección criptográfica de credenciales.
Su implementación permite estudiar cómo interactúan estos elementos en una aplicación de autenticación de alcance local y cómo una interfaz textual puede proporcionar navegación, mensajes de estado y formularios estructurados sin recurrir a un entorno gráfico.

\noindent\textbf{Palabras clave:} sistemas operativos, lenguaje C, \texttt{ncurses}, autenticación, Argon2id, persistencia local y validación de datos.

\newpage

% ÍNDICE
\tableofcontents

\newpage

% INTRODUCCIÓN
\section{Introducción}
Los sistemas de inicio de sesión constituyen un mecanismo fundamental para controlar el acceso a las aplicaciones que administran información asociada con distintos usuarios.
Su funcionamiento requiere identificar a quien intenta ingresar, verificar sus credenciales y gestionar adecuadamente los datos utilizados durante el registro y la autenticación.
Además de comprobar la correspondencia entre un nombre de usuario y una contraseña, resulta necesario contemplar la validación de las entradas, el almacenamiento persistente de las cuentas y la presentación de mensajes que permitan comprender el resultado de cada operación.

En el contexto de los sistemas operativos, el desarrollo de una aplicación de estas características permite estudiar la interacción entre un programa y el entorno en el que se ejecuta.
El manejo del teclado, la presentación de información en una terminal, la lectura y escritura de archivos y el uso de bibliotecas externas forman parte de los elementos que intervienen en su funcionamiento.
Estos componentes adquieren especial relevancia cuando la aplicación debe mantener una interacción organizada sin utilizar ventanas, botones gráficos ni otros elementos propios de una interfaz visual tradicional.

A partir de este contexto, se desarrolló JDM E-TICKETS, un sistema de registro e inicio de sesión implementado en lenguaje C.
La aplicación utiliza la biblioteca \texttt{ncurses} para construir una interfaz basada en texto que permite navegar por las opciones y completar formularios mediante el teclado.
Su funcionamiento contempla dos operaciones principales: la creación de cuentas nuevas y la autenticación de usuarios registrados.
Para el registro, se establecen condiciones sobre el nombre de usuario, la dirección de correo electrónico y la contraseña, mientras que el inicio de sesión se basa en la existencia de la cuenta y la verificación de la contraseña proporcionada.

Otro aspecto relevante del proyecto es el tratamiento de las credenciales.
En lugar de conservar las contraseñas originales, la aplicación emplea \texttt{libsodium} y el algoritmo Argon2id para generar hashes que posteriormente se utilizan durante la autenticación.
Los registros se almacenan en un archivo local denominado \texttt{usuarios.txt}, lo que permite conservar la información entre ejecuciones sin incorporar una base de datos ni depender de un servidor externo.
Esta decisión delimita el alcance de la solución y permite concentrar su implementación en el procesamiento de datos, la interacción con el usuario y las operaciones de archivos.

El objetivo principal del proyecto es implementar un mecanismo funcional de registro y autenticación que integre una interfaz textual, validaciones de entrada, verificación criptográfica de contraseñas y persistencia local.
Para ello, el código se organiza en módulos que separan la lógica principal, las funciones de autenticación y los componentes visuales, mientras que el proceso de compilación se administra mediante un archivo \texttt{Makefile}.

El reporte presenta los fundamentos relacionados con el funcionamiento del sistema y describe las decisiones tomadas durante su implementación.
Asimismo, aborda la estructura del código fuente, la compilación, la construcción de la interfaz, las reglas de validación, el tratamiento de las contraseñas y el almacenamiento de los usuarios.
Finalmente, se incluyen las conclusiones correspondientes y las referencias utilizadas para contextualizar el desarrollo.
De esta manera, se expone la relación entre los conceptos estudiados y su aplicación en un programa de autenticación ejecutado desde una terminal de Linux.

\newpage

% IMPORTACIÓN DE MÓDULOS
\section{Marco Teórico}

\subsection{Teoría sobre la creación de un inicio de sesión}
1.1
La autenticación de usuarios es el mecanismo fundamental que permite a un sistema operativo validar la identidad de los sujetos que solicitan acceso a sus recursos. De acuerdo con Silberschatz, Galvin y Gagne en Fundamentos de Sistemas Operativos, es indispensable establecer una distinción clara entre la identificación y la autenticación. La identificación ocurre cuando un usuario declara su identidad mediante una credencial o nombre de cuenta. Por su parte, la autenticación comprende la comprobación explícita de dicha declaración. Las técnicas de autenticación se dividen fundamentalmente en tres categorías: información que el usuario conoce, objetos que el usuario posee e indicadores biométricos que identifican sus características físicas o conductuales.

En el contexto de la administración de accesos en entornos tipo UNIX, Francisco Manuel Márquez García describe en Unix Programación Avanzada la secuencia técnica de llamadas al sistema encargadas del flujo de inicio de sesión. La arquitectura de estos sistemas delega la gestión de las líneas de terminal al proceso inicializador del sistema, el cual invoca rutinas especializadas para configurar la línea de comunicación y solicitar el identificador del usuario. Posteriormente, se transfiere el control al ejecutable de inicio de sesión, el cual desactiva el eco en la terminal, solicita la clave de acceso y consulta los archivos de seguridad del sistema antes de inicializar el entorno mediante el establecimiento de identificadores de usuario y la ejecución del intérprete de comandos.

La arquitectura de protección de credenciales en el sistema operativo exige una separación rigurosa en el almacenamiento de datos. Como exponen Andrew S. Tanenbaum y Herbert Bos en Sistemas Operativos Modernos, las contraseñas nunca deben almacenarse en texto plano. Históricamente, la información técnica de los usuarios se concentraba en el archivo del sistema de acceso público; sin embargo, los sistemas modernos implementan mecanismos de protección que restringen los permisos de lectura únicamente al usuario con privilegios de administrador, albergando los resúmenes cifrados de las contraseñas junto con los parámetros de expiración de las cuentas.

1.1.1. Identificación y autenticación de usuarios
El proceso de identificación constituye el primer paso dentro del control de acceso en un sistema operativo multiusuario. En esta etapa, el usuario proporciona un identificador único, como un nombre de usuario o un número de cuenta, para declararse ante el sistema. No obstante, la sola declaración de la identidad es insuficiente para otorgar privilegios sobre los recursos del sistema operativo, por lo que se requiere un procedimiento riguroso de autenticación. Silberschatz, Galvin y Gagne destacan que la autenticación verifica la validez de la afirmación presentada. La forma más extendida de autenticación se basa en secretos compartidos, como las contraseñas numéricas o alfanuméricas, donde el usuario demuestra su conocimiento sobre un dato privado previamente registrado. En entornos que requieren mayor grado de seguridad, los sistemas operativos integran múltiples factores combinados para reducir las vulnerabilidades derivadas del robo o la adivinación de claves de acceso.

1.1.2. El flujo de inicio de sesión en sistemas UNIX
En los sistemas operativos basados en la filosofía UNIX, el inicio de sesión es administrado mediante una jerarquía de procesos encadenados. Márquez García explica que el proceso raíz o inicializador del sistema genera procesos dedicados a supervisar cada terminal o línea de entrada disponible. Cuando un usuario interactúa con la terminal, la rutina getty establece la velocidad de transmisión, los caracteres de control de la línea y despliega el mensaje de bienvenida solicitando el nombre de usuario. Tras la recepción del nombre de usuario, dicho proceso sustituye su imagen en memoria mediante la llamada al sistema execve para invocar al programa login. El ejecutable login se encarga de cambiar la configuración del controlador de la terminal para deshabilitar el eco del teclado, solicita la contraseña, aplica el algoritmo de verificación y, en caso de autenticación exitosa, asigna los identificadores de usuario y grupo correspondientes mediante llamadas al sistema como setuid y setgid, finalizando con la ejecución del shell configurado para el usuario.

1.1.3. Archivos de control y protección de credenciales (/etc/passwd y /etc/shadow)
La estructura de almacenamiento de credenciales en sistemas UNIX evolucionó debido a las limitaciones de seguridad del diseño original. Tanenbaum y Bos detallan que tradicionalmente el archivo /etc/passwd guardaba todos los datos de la cuenta de usuario, incluyendo el nombre, la contraseña cifrada, el identificador de usuario, el identificador de grupo predeterminado, el directorio de trabajo y el intérprete de comandos asignado. Dado que numerosos programas legítimos del sistema necesitan consultar la información de los usuarios en /etc/passwd, dicho archivo debe conservar permisos de lectura global para todos los procesos del sistema. Esta visibilidad pública permitía que usuarios malintencionados copiaran el archivo y aplicaran técnicas de fuerza bruta o diccionarios fuera de línea sobre las contraseñas cifradas. Para erradicar este riesgo, se introdujo la técnica del archivo en sombra o /etc/shadow, el cual separa las contraseñas cifradas y las traslada a un archivo protegido accesible únicamente por el usuario administrador. En /etc/passwd el campo de la contraseña se sustituye por una marca fija, mientras que /etc/shadow resguarda los datos sensibles, la sal de cifrado y los parámetros de expiración de la clave.


\subsection{Algoritmos de cifrado}
El almacenamiento de contraseñas en texto plano representa una vulnerabilidad crítica. Para mitigar este riesgo, los sistemas implementan algoritmos de cifrado o funciones de un solo sentido (hashing). En lenguaje C, es común utilizar funciones matemáticas o bibliotecas nativas como \texttt{crypt()} (basada en DES, MD5 o SHA) que transforman la contraseña original en una cadena de caracteres ilegible. Durante el inicio de sesión, el sistema cifra la contraseña ingresada y compara el resultado con el hash guardado; si ambos coinciden, se otorga el acceso sin que el sistema necesite conocer la clave original.

\subsection{Biblioteca ncurses}
\texttt{ncurses} (new curses) es una biblioteca de programación que proporciona una API para construir interfaces de usuario basadas en texto de manera independiente de la terminal. Permite controlar la posición del cursor, crear ventanas, utilizar colores y leer el teclado carácter por carácter. Funciones como \texttt{noecho()} son esenciales en los inicios de sesión, ya que desactivan el eco visual de la terminal, permitiendo que el usuario escriba su contraseña sin que los caracteres sean visibles en la pantalla.

\newpage

\section{Desarrollo}

\subsection{Reglas de negocio y Compilación (Makefile)}
Para esta práctica, todo el código fuente distribuido en los archivos \texttt{.c} y \texttt{.h} fue estructurado para compilarse mediante un archivo \texttt{Makefile}. A continuación se demuestra que la compilación se ejecutó estrictamente en lenguaje C sin advertencias ni errores:

% Descomentar y ajustar la ruta de la imagen cuando tengan la captura
% \begin{figure}[H]
%     \centering
%     \includegraphics[width=0.7\textwidth]{img/captura_makefile.png}
%     \caption{Compilación limpia del código fuente mediante el comando make.}
% \end{figure}

\subsection{Implementación del inicio de sesión con ncurses}
Se implementó un menú interactivo utilizando la biblioteca ncurses para estructurar visualmente la entrada de datos. El sistema solicita el nombre de usuario y la contraseña, habiendo configurado previamente la terminal con la función \texttt{noecho()} para ocultar el texto introducido en el campo de la clave, cumpliendo con las especificaciones de la práctica.

% Descomentar y ajustar la ruta de la imagen cuando tengan la captura
% \begin{figure}[H]
%     \centering
%     \includegraphics[width=0.7\textwidth]{img/captura_login.png}
%     \caption{Interfaz basada en texto (ncurses) solicitando usuario y contraseña oculta.}
% \end{figure}

\subsection{Cifrado de contraseñas y validación}
Antes de procesar las credenciales, el sistema valida que los campos no se encuentren vacíos y no excedan el límite de caracteres permitidos para prevenir desbordamientos de búfer. Seguidamente, la contraseña en texto plano pasa por el algoritmo de cifrado antes de ser comparada con los registros existentes o guardada en caso de ser un usuario nuevo. No se utilizan gestores de bases de datos externos.

\subsection{Registro y persistencia en archivos}
Los usuarios y sus respectivas contraseñas cifradas se administran utilizando descriptores de archivo nativos de C. Si el sistema detecta que el usuario no se encuentra registrado, despliega la opción de registro. El sistema abre el archivo en modo de adición (\texttt{a+}) y escribe la nueva estructura de credenciales, garantizando la persistencia de la información para ejecuciones futuras.

% Descomentar y ajustar la ruta de la imagen cuando tengan la captura
% \begin{figure}[H]
%     \centering
%     \includegraphics[width=0.7\textwidth]{img/captura_registro.png}
%     \caption{Formulario de registro para nuevos usuarios y almacenamiento en el sistema de archivos.}
% \end{figure}

\newpage
\FloatBarrier
\section{Conclusiones}

\subsection{Alvarez Garcia Martin}
% Escribe tu conclusión aquí. Recuerda usar una oración por línea de código para evitar conflictos en Git.

\vspace{1cm}

\subsection{Hernandez Castillo Emir Adrian}
% Escribe tu conclusión aquí. Recuerda usar una oración por línea de código para evitar conflictos en Git.

\vspace{1cm}

\subsection{Joaquín Rojas González}
% Escribe tu conclusión aquí. Recuerda usar una oración por línea de código para evitar conflictos en Git.

\vspace{1cm}

\subsection{Nava Alvarado Donovan Emilio}
% Escribe tu conclusión aquí. Recuerda usar una oración por línea de código para evitar conflictos en Git.

\newpage

% REFERENCIAS
\newpage
\addcontentsline{toc}{section}{Referencias}
\section{Referencias}

\begin{thebibliography}{00}

\bibitem{tanenbaum2009}
A. S. Tanenbaum, \textit{Sistemas Operativos Modernos}, 3.a ed.
México: Pearson Educación, 2009.

\bibitem{silberschatz2006}
A. Silberschatz, P. B. Galvin y G. Gagne,
\textit{Fundamentos de Sistemas Operativos}, 7.a ed.
México: McGraw Hill Education, 2006.

\bibitem{stevens2013}
W. R. Stevens y S. A. Rago,
\textit{Advanced Programming in the UNIX Environment}, 3.a ed.
Addison-Wesley Professional, 2013.

\bibitem{kerrisk2010}
M. Kerrisk,
\textit{The Linux Programming Interface: A Linux and UNIX System Programming Handbook}.
San Francisco, CA: No Starch Press, 2010.

\bibitem{libsodium}
Libsodium. \textit{Password Hashing}. Documentación oficial.
Disponible en: \url{https://doc.libsodium.org/password_hashing}.

\bibitem{make}
GNU Project. \textit{GNU Make Manual}.
Disponible en: \url{https://www.gnu.org/software/make/manual/make.html}.

\end{thebibliography}

\newpage

% ANEXOS
\section*{Anexos}
\phantomsection
\addcontentsline{toc}{section}{Anexos}
En esta sección se incluye el repositorio oficial donde se encuentra el código fuente en lenguaje C, los archivos de cabecera y el archivo Makefile necesarios para la compilación del sistema:
\begin{center}
    \url{https://github.com/StrlgE26/p3so.git}
\end{center}

\end{document}